\section{Experiments}

\subsection{Datasets}

To comprehensively evaluate the proposed Unified Dual-Mask Physical Model and validate its generalization across varying surface reflectance properties, we construct a multi-domain experimental setup using three representative light field (LF) datasets. These datasets are meticulously selected to cover a wide spectrum of scene types, ranging from ideal Lambertian surfaces to highly complex mixed urban environments and physically challenging non-Lambertian materials. 

\textbf{HCInew Dataset.} The HCInew dataset primarily serves as the benchmark for \textbf{Lambertian} and standard synthetic scenes. It comprises 22 high-quality scenes (14 scenes for training and 8 for testing) (e.g., \textit{boxes}, \textit{cotton}, \textit{dino}) rendered using Blender [4]. Each scene consists of $9 \times 9$ (81) angular views with a spatial resolution of $512 \times 512$ pixels. The ground-truth (GT) depth maps are perfectly aligned and generated directly from the 3D rendering engine, providing pixel-accurate depth supervision. The dataset is divided into training and testing sets following the standard stratified split (e.g., `training/stratified`), ensuring a diverse distribution of geometric structures in both splits. It is chosen to evaluate the model's fundamental capability in extracting Epipolar Plane Image (EPI) structures under ideal physical assumptions.

\textbf{UrbanLF-Syn Dataset.} To assess the model's robustness in \textbf{Mixed/Urban} scenarios, we incorporate the UrbanLF-Syn dataset . This dataset contains 170 synthetic scenes depicting complex urban street views, featuring intricate textures, severe occlusions, and large-scale architectural structures. Similar to HCInew, it provides $9 \times 9$ angular views with high spatial resolution. The GT depth maps are derived from high-fidelity 3D city models. The scenes are partitioned into training and testing subsets. We select this dataset to challenge the model with complex texture variations and real-world structural complexities, where EPI slopes might be locally ambiguous due to fine details and occlusions.

\textbf{Non-Lambertian Dataset.} Evaluating depth estimation on \textbf{Non-Lambertian} surfaces is the core focus of this work. We utilize a specialized Non-Lambertian dataset (derived from the Zhenglong/Wanner HCI collections ) that explicitly features specular (highlight), translucent, and scattering surfaces. Each scene contains $9 \times 9$ angular views at a $512 \times 512$ resolution, with GT depth maps generated via synthetic rendering. \textit{Preprocessing and Constraints:} Originally, the HCI-Old dataset was considered; however, it was excluded from our pipeline due to inconsistent data formatting and lack of standardized depth calibration parameters. Consequently, the available Non-Lambertian dataset is severely constrained, containing only 4 training scenes. Despite this extreme data scarcity, it is indispensable for our study as it introduces the ultimate physical challenge: the violation of the standard EPI linearity assumption caused by view-dependent reflectance (e.g., specular highlights shifting across views).

\textbf{Data Preprocessing and Domain-Balanced Sampling.} 
A critical challenge in our multi-domain setup is the severe data imbalance across different scene types (e.g., 170 scenes in UrbanLF-Syn versus only 4 scenes in the Non-Lambertian dataset). To prevent the network from overfitting to the data-rich domains and neglecting the physically challenging non-Lambertian cases, we implement a rigorous \textbf{domain-balanced sampling strategy}. Specifically, we employ a `WeightedRandomSampler` during the training phase. Each sample is assigned a weight inversely proportional to the total number of samples in its respective domain (Lambertian, Mixed/Urban, or Non-Lambertian). This ensures that each mini-batch maintains an equitable exposure to all three domains, facilitating stable multi-task learning and joint optimization of the depth and domain classification losses.

The key characteristics and configurations of the datasets utilized in our experiments are summarized in Table I.

\textbf{TABLE I}
\textbf{SUMMARY OF DATASETS USED FOR TRAINING AND EVALUATION}

\begin{table*}[!t]
\small
\caption{Dataset Domain / Scene Type No. of Scenes (Train / Test).}
\label{tab:table1}
\centering
\resizebox{\textwidth}{!}{
\begin{tabular*}{\textwidth}{@{\extracolsep{\fill}}lllllll}
\toprule
Dataset \& Domain / Scene Type \& No. of Scenes (Train / Test) \& Angular Views \& Spatial Resolution \& GT Depth Source \& Key Characteristics \& Challenges \\
\midrule
\textbf{HCInew} \& Lambertian \& 22 (14 / 8) \& $9 \times 9$ (81) \& $512 \times 512$ \& Synthetic (Blender) \& Ideal EPI structures; standard geometric shapes and textures. \\
\textbf{UrbanLF-Syn} \& Mixed / Urban \& 170 (Partitioned) \& $9 \times 9$ (81) \& $512 \times 512$ \& Synthetic (3D City Models) \& Complex textures, severe occlusions, large-scale urban structures. \\
\textbf{Non-Lambertian} \& Non-Lambertian \& 4 / (Cross-validated) \& $9 \times 9$ (81) \& $512 \times 512$ \& Synthetic (Blender) \& Specular, translucent, scattering surfaces; EPI physical assumption failure; extreme data scarcity. \\
\bottomrule
\end{tabular*}
}
\end{table*}
\textit{Note: The HCI-Old dataset was excluded during preprocessing due to inconsistent data formatting and lack of standardized depth calibration parameters. Consequently, the Non-Lambertian training set is limited to 4 scenes.}

\subsubsection{Implementation Details}

\textbf{Hardware and Software Environment.} 
The proposed Unified Dual-Mask Physical Model is implemented using the PyTorch framework. All experiments are conducted on a high-performance workstation equipped with four NVIDIA RTX 3090 GPUs (24GB VRAM each) and an AMD EPYC 7742 CPU. The Distributed Data Parallel (DDP) strategy is utilized to synchronize gradients and accelerate the training process across multiple GPUs.

\textbf{Training Hyperparameters.} 
We employ the AdamW optimizer to train the network, initialized with a learning rate of $1 \times 10^{-4}$ and a weight decay of $1 \times 10^{-4}$. The learning rate is dynamically adjusted using a Cosine Annealing scheduler, decaying to a minimum of $1 \times 10^{-6}$ over the course of training. The batch size is set to 32 per GPU, and gradient accumulation is applied with 2 steps, yielding an effective global batch size of 256. The model is trained for a total of 150 epochs. To mitigate the severe data imbalance among the Lambertian, Mixed/Urban, and Non-Lambertian datasets (notably, the Non-Lambertian dataset contains only 4 training scenes), we adopt a domain-balanced sampling strategy utilizing a `WeightedRandomSampler`. This ensures equitable exposure to each physical domain during every epoch. The key hyperparameters are summarized in Table II.

\textbf{TABLE II}
\textbf{SUMMARY OF KEY TRAINING HYPERPARAMETERS}

\begin{table}[!t]
\caption{Hyperparameter Value.}
\label{tab:table2}
\centering
\begin{tabular}{ll}
\toprule
Hyperparameter \& Value \\
\midrule
Optimizer \& AdamW \\
Initial Learning Rate \& $1 \times 10^{-4}$ \\
LR Scheduler \& Cosine Annealing \\
Minimum Learning Rate \& $1 \times 10^{-6}$ \\
Batch Size (per GPU) \& 32 \\
Gradient Accumulation Steps \& 2 \\
Effective Batch Size \& 256 \\
Total Epochs \& 150 \\
Weight Decay \& $1 \times 10^{-4}$ \\
Sampler \& WeightedRandomSampler (Domain-balanced) \\
\bottomrule
\end{tabular}
\end{table}
\textbf{Data Augmentation.} 
To improve the generalization capability and prevent overfitting—particularly given the extreme scarcity of Non-Lambertian training samples—we apply a comprehensive set of data augmentation techniques tailored for 4D light field data. These include random spatial cropping (resizing to $512 \times 512$ resolution), random horizontal and vertical flips (applied consistently across all angular views to preserve epipolar geometry), and random angular shifts. Additionally, we apply color jittering (random adjustments to brightness, contrast, and saturation) and Gaussian noise injection. These photometric augmentations simulate sensor noise and varying illumination conditions, which is highly beneficial for robust material recognition and domain classification.

\textbf{Loss Function Weights.} 
The network is optimized in a multi-task learning paradigm, jointly minimizing the depth estimation loss and the auxiliary domain classification loss. The overall objective function is defined as $\mathcal{L}_{total} = \lambda_{depth} \mathcal{L}_{depth} + \lambda_{domain} \mathcal{L}_{domain}$. We empirically set the loss weights to $\lambda_{depth} = 1.0$ and $\lambda_{domain} = 0.1$ to prioritize accurate depth regression while providing sufficient gradient signals for domain-aware feature disentanglement.

\subsubsection{Evaluation Metrics}

To quantitatively assess the performance of the proposed Unified Dual-Mask Physical Model and ensure fair comparisons with existing state-of-the-art methods, we employ a comprehensive set of evaluation metrics. Following standard protocols, we use Mean Squared Error (MSE) and Bad Pixel ratio as our primary evaluation metrics for depth estimation accuracy.

Specifically, the \textbf{Mean Squared Error (MSE)} measures the average squared difference between the estimated depth map and the ground-truth depth map, providing a global assessment of regression accuracy. Lower MSE values indicate better overall depth estimation performance.

Furthermore, to evaluate local precision and the proportion of highly inaccurate predictions, we compute the \textbf{Bad Pixel ratio}. A pixel is classified as a "bad pixel" if the absolute error between the predicted depth and the ground truth exceeds a predefined threshold. Following standard protocols in light field depth estimation, we report the Bad Pixel ratios at two strict thresholds: $<0.01$ and $<0.03$ (normalized by the maximum depth range). A lower Bad Pixel ratio signifies that the model produces fewer severe outliers, which is particularly crucial for challenging non-Lambertian regions where EPI structures are heavily distorted.

Additionally, we report the \textbf{Mean Absolute Error (MAE)} to provide a robust measure of the average magnitude of errors, complementing the MSE metric. All metrics are computed exclusively on the valid regions of the depth maps, ignoring masked or invalid background pixels.